% =====================================================================
%  METHODOLOGY — NeuroFlow
%  Journal-style replacement section. Few subsections, no code refs.
%  Covers every step required by the current Results:
%    - VENC-normalised intraluminal MSE (total + per component)
%    - Background MSE
%    - Geometry metrics: Dice, sMSD, HD
%    - Two input strategies: synthetic-noise and masked-input
%    - Leave-One-Out Cross-Validation (n=7)
%  Required preamble packages: amsmath, amssymb, graphicx, booktabs.
% =====================================================================

\section*{Methodology}

\subsection*{Study population and 4D Flow MRI acquisition}

Seven patients with cerebrovascular pathology were prospectively
enrolled and scanned on two MRI platforms within a single imaging
session: a clinical 3\,T Siemens system and an investigational 7\,T
Siemens system. For each subject, paired 4D Flow MRI studies were
acquired on both scanners to minimise inter-session physiological
variability. Each acquisition produced four co-registered volumetric
time series --- a magnitude image and three directional phase-contrast
velocity maps encoding velocity along the feet--head ($u$),
left--right ($v$), and anterior--posterior ($w$) axes
\cite{markl2012,dyverfeldt2015}. The 3\,T series covered between nine
and fourteen cardiac frames, while the 7\,T series spanned ten to
fifteen frames, all encoded with a common velocity-encoding sensitivity
of $\mathrm{VENC} = 0.9$\,\textnormal{m/s} across both field strengths
\cite{bock2010}. Raw data were converted from DICOM to NIfTI, preserving
the original phase representation and applying the LPS\,$\rightarrow$\,RAS
sign convention so that the three velocity components retained their
physical orientation in the scanner frame.

% FIGURE — Acquisition and pairing overview.
% Place a multi-panel figure illustrating the dual-field-strength
% acquisition protocol: representative axial slices of the paired 3T and
% 7T magnitude (same anatomical location), sagittal MIPs of each
% acquisition to convey 3D coverage, and a schematic of the within-session
% pairing.
\begin{figure}[!htbp]
    \centering
    % Insert figure here
    \caption{Acquisition and pairing overview. Paired 3\,T and 7\,T 4D
    Flow MRI acquired within a single imaging session for each subject.}
    \label{fig:acquisition}
\end{figure}

\subsection*{Motion correction and multi-field-strength registration}

Cardiac-gated 4D Flow MRI is intrinsically sensitive to bulk motion
between cardiac phases. Intra-scan motion was corrected by estimating
rigid-body transformations $\{\mathbf{T}_t^{\text{rigid}}\}_{t=1}^{T}$
from the magnitude time series relative to a reference frame using
normalised mutual information \cite{maes1997}, and applying each
transformation consistently to the paired velocity volume,
\begin{equation}
    \mathbf{V}_t^{\text{corr}}
    = \mathbf{T}_t^{\text{rigid}} \circ \mathbf{V}_t^{\text{raw}},
    \qquad t = 1, \ldots, T,
\end{equation}
followed by a vector-aware reorientation of the velocity components. To
bring the 7\,T volumes into the geometric space of the 3\,T
acquisitions, a temporal maximum-intensity projection
$\mathrm{MIP}(\mathbf{x}) = \max_t M_t(\mathbf{x})$ was first computed
for each magnitude series to obtain a static, high-contrast angiographic
representation. A rigid alignment followed by an affine refinement,
both optimised with normalised mutual information, defined the
composite transformation $\mathbf{T}^{\text{affine}}_{7T \to 3T}$, which
was then applied to the full 4D velocity volume so that
\begin{equation}
    \mathbf{V}_{7T \to 3T}^{\text{reg}}(\mathbf{x}, t)
    = \mathbf{T}^{\text{affine}}_{7T \to 3T} \circ
      \mathbf{V}_{7T}^{\text{corr}}(\mathbf{x}, t),
\end{equation}
establishing voxel-level correspondence between the paired 3\,T and
7\,T data as required for supervised learning.

% FIGURE — Registration pipeline.
% Multi-panel figure: (a) temporal MIP of a 3T magnitude image; (b)
% temporal MIP of the corresponding 7T magnitude image before
% registration; (c) 7T MIP after rigid + affine registration into 3T
% space, with a checkerboard or difference overlay to demonstrate
% alignment quality; (d) optional displacement-magnitude field.
\begin{figure}[!htbp]
    \centering
    % Insert figure here
    \caption{Registration pipeline. Temporal MIPs of the paired 3\,T and
    7\,T magnitude before and after the rigid + affine inter-scan
    registration; the checkerboard overlay illustrates the resulting
    voxel-level alignment.}
    \label{fig:registration}
\end{figure}

\subsection*{Angiographic volume synthesis and vascular masking}

A voxel-level vascular probability map was constructed by combining
two complementary contrast sources. A 4D speed image was obtained from
the VENC-scaled velocity vector,
\begin{equation}
    S(\mathbf{x}, t) = \sqrt{
        \left(\tfrac{u(\mathbf{x},t)}{\mathrm{VENC}}\right)^2 +
        \left(\tfrac{v(\mathbf{x},t)}{\mathrm{VENC}}\right)^2 +
        \left(\tfrac{w(\mathbf{x},t)}{\mathrm{VENC}}\right)^2},
\end{equation}
and projected over the cardiac cycle as
$\bar{S}(\mathbf{x}) = \max_t S(\mathbf{x},t)$. The temporal MIP of the
magnitude and $\bar{S}$ were each min--max normalised and their
voxel-wise product
$\mathcal{A}(\mathbf{x}) = \hat{\mathrm{MIP}}(\mathbf{x}) \cdot
\hat{S}(\mathbf{x})$ defined a synthetic angiographic volume in which
voxels that are both bright in magnitude and fast in flow are reinforced
while static tissue and background noise are suppressed.

The Circle-of-Willis (CoW) mask used throughout this work was generated
on the 7\,T magnitude data, taking advantage of their substantially
higher signal-to-noise ratio. For every subject, an initial foreground
estimate was obtained by adaptive intensity thresholding of the 7\,T
maximum-intensity projection with connected-component cleanup, then
refined by an ensemble nnU-Net inference and finally corrected by an
expert reviewer on a slice-by-slice basis. The resulting manually
corrected mask defined the gold-standard CoW segmentation in the 7\,T
geometry; transferring it to the 3\,T frame via the registration above
yielded the binary mask $\mathcal{M}(\mathbf{x})$ used as the
supervision and evaluation region.

Two nnU-Net models with the residual-encoder \emph{ResEnc-M} backbone
\cite{isensee2021nnunet} were then trained on these corrected labels
with the skeleton-recall loss formulation introduced by the
TopCoW-2024 winning solution \cite{yang2023topcow,aydin2024topcow}.
The objective combined Dice, skeleton-recall, and cross-entropy terms,
\begin{equation}
    \mathcal{L}_{\text{seg}}
    = w_{D}\,\mathcal{L}_{\text{Dice}}
    + w_{S}\,\mathcal{L}_{\text{SkelRec}}
    + w_{C}\,\mathcal{L}_{\text{CE}},
    \qquad w_{D}=w_{S}=w_{C}=1,
\end{equation}
where the skeleton-recall term is computed on the fly from a
differentiable topological skeletonisation of the ground-truth mask.
Both 2D and 3D-full-resolution configurations were trained over five
cross-validation folds and combined as a 2D\,$+$\,3D ensemble at
inference. The first model operates on the temporal maximum projection
of the native 3\,T magnitude (mean cross-validated Dice~$0.83$), and a
second model was trained on the same projection after isotropic
$0.5{\times}$ downsampling (mean Dice~$0.76$), matching the input
geometry of the super-resolution task introduced below. The same
ensemble is used twice in this work: first, at training time, to
propagate the 7\,T-derived gold-standard mask through the 3\,T,
$0.5{\times}$\,3\,T, and 7\,T-in-3\,T volumes that form the supervised
pairs; and second, at evaluation time, to re-segment each predicted
4D Flow volume so that its CoW geometry can be compared against the
subject-specific 7\,T reference mask.

% FIGURE — Angiographic synthesis and vascular masking.
% Two-row composite: top row shows axial and coronal views of (a)
% normalised temporal MIP of the magnitude, (b) normalised temporal speed
% projection, and (c) the synthesised angiographic product A(x); bottom
% row shows the ground-truth provenance pipeline on a single 7T MIP slice
% (hybrid threshold -> bootstrap nnU-Net -> manually corrected mask) and
% a small donut chart with the 5-fold cross-validated Dice for both
% segmentation models (3T and 0.5 x 3T).
\begin{figure}[!htbp]
    \centering
    % Insert figure here
    \caption{Angiographic volume synthesis and vascular masking. The
    synthesised angiographic field $\mathcal{A}(\mathbf{x})$ combines
    magnitude and speed evidence; the gold-standard CoW mask is derived
    from the 7\,T magnitude via a hybrid threshold, bootstrap nnU-Net,
    and manual correction, and a ResEnc-M U-Net trained with the
    skeleton-recall loss propagates this mask to the 3\,T and
    $0.5{\times}3$\,T inputs and to the predicted volumes used in
    evaluation.}
    \label{fig:segmentation}
\end{figure}

\subsection*{Network architecture: CBAM-DNS and CBAM-SupRes}

The core predictive backbone is an adapted SR4DFlowNet
\cite{ferdian2020,elahmar2023}, a 3D fully convolutional residual
network inspired by enhanced residual architectures for super-resolution
\cite{lim2017edsr}, augmented with a Convolutional Block Attention
Module (CBAM) \cite{woo2018cbam} inserted between the low-resolution
residual stack and the upsampling stage. We instantiate the same
backbone in two task-specific forms used throughout the results:
\textbf{CBAM-DNS} for same-resolution 3\,T$\to$7\,T domain translation
($r=1$) and \textbf{CBAM-SupRes} for joint domain translation and
twofold super-resolution ($r=2$).

The model accepts six scalar input channels per voxel --- the three
VENC-normalised velocity components $(u, v, w)$ and their
magnitude-weighted counterparts $(u_m, v_m, w_m)$ --- and internally
derives three auxiliary fields that summarise the local kinematics and
phase-contrast signal:
\begin{equation}
    S = \sqrt{u^{2}+v^{2}+w^{2}},
    \qquad
    M = \sqrt{u_m^{2}+v_m^{2}+w_m^{2}},
    \qquad
    \mathrm{PCMR} = M \cdot S.
\end{equation}
Velocity and physical channels are processed in two dedicated
convolutional streams of two $3{\times}3{\times}3$ layers with 64
feature maps, symmetric padding, and ReLU activation: a phase stream
operating on $(u,v,w)$ to extract flow-kinematic features, and a
phase-contrast stream operating on $(\mathrm{PCMR}, M, S)$ to extract
signal-intensity features. Their feature maps are concatenated and
fused by a $1{\times}1{\times}1$ projection followed by a
$3{\times}3{\times}3$ refinement convolution.

The fused representation is propagated through $N_{LR}=8$ pre-activation
residual blocks at the input resolution \cite{he2016resnet},
\begin{equation}
    \mathbf{F}^{(l+1)} = \mathbf{F}^{(l)}
        + f\!\left(\mathbf{F}^{(l)};\, \Theta^{(l)}\right),
\end{equation}
after which the CBAM module sequentially applies channel and spatial
attention to re-weight feature responses before upsampling. Channel
attention is obtained from global average- and max-pooled descriptors
processed by a shared two-layer multilayer perceptron, while spatial
attention is computed from the channel-wise pooled feature maps and a
$7{\times}7{\times}7$ convolution; the two attention maps are applied
multiplicatively in sequence. Spatial upsampling is then performed by
trilinear interpolation with scale factor $r$, followed by $N_{HR}=4$
high-resolution residual blocks operating at the target resolution. The
network terminates in three independent two-layer convolutional heads
that emit the reconstructed velocity components
$(\hat{u}, \hat{v}, \hat{w})$, together with an additional head that
predicts the reconstructed magnitude $\hat{M}$. Within this common
backbone, CBAM-DNS uses $r=1$ and trains under the same-resolution
domain-translation setting, whereas CBAM-SupRes uses $r=2$ and absorbs
both the cross-field-strength translation and the spatial upsampling
into a single forward pass.

% FIGURE — Network architecture diagram.
% Full architectural diagram of CBAM-DNS / CBAM-SupRes: dual-stream
% input (phase and PC paths), feature concatenation and 1x1 fusion, the
% low-resolution residual stack, the CBAM channel + spatial attention
% block, the trilinear upsampling block (r = 1 or 2), the high-resolution
% residual stack, and the velocity and magnitude output heads. Inset
% showing the CBAM module structure.
\begin{figure}[!htbp]
    \centering
    % Insert figure here
    \caption{Network architecture underlying CBAM-DNS ($r=1$) and
    CBAM-SupRes ($r=2$). The phase and phase-contrast streams are fused,
    refined by $N_{LR}=8$ residual blocks, re-weighted by a CBAM
    channel-and-spatial attention module, upsampled by trilinear
    interpolation with factor $r$, and refined again by $N_{HR}=4$
    high-resolution residual blocks before three velocity heads and an
    optional magnitude head emit the reconstructed volume.}
    \label{fig:architecture}
\end{figure}

\subsection*{Loss function and evaluation metrics}

Training minimises a composite, mask-aware loss that jointly penalises
intraluminal velocity error, magnitude error, and texture irregularity
in the extraluminal background,
\begin{equation}
    \mathcal{L} = \mathcal{L}_{\text{vel}}
                + \lambda_{M}\,\mathcal{L}_{\text{mag}}
                + \lambda_{\text{TV}}\,\mathcal{L}_{\text{TV}}.
\end{equation}
For each predicted velocity component
$\hat{q}\in\{\hat{u},\hat{v},\hat{w}\}$ and its VENC-normalised target
$q^{*}=q/\mathrm{VENC}$, the per-voxel squared error
$\epsilon_q(\mathbf{x})=(\hat{q}(\mathbf{x})-q^{*}(\mathbf{x}))^{2}$ is
split into a masked fluid term and a complementary background term,
\begin{equation}
    \mathcal{L}^{q}_{\text{fluid}}
    = \frac{\sum_{\mathbf{x}} \epsilon_q(\mathbf{x})\,\mathcal{M}(\mathbf{x})}
           {\sum_{\mathbf{x}} \mathcal{M}(\mathbf{x}) + \varepsilon},
    \qquad
    \mathcal{L}^{q}_{\text{bg}}
    = \frac{\sum_{\mathbf{x}} \epsilon_q(\mathbf{x})\,(1-\mathcal{M}(\mathbf{x}))}
           {\sum_{\mathbf{x}} (1-\mathcal{M}(\mathbf{x})) + \varepsilon},
\end{equation}
and aggregated as
$\mathcal{L}_{\text{vel}}
 = \sum_{q\in\{u,v,w\}}
   (\mathcal{L}^{q}_{\text{fluid}}+\lambda_{\text{bg}}\,\mathcal{L}^{q}_{\text{bg}})$
with background weight $\lambda_{\text{bg}}=0.1$, reflecting the
clinical priority of accurate intraluminal velocity quantification.
$\mathcal{L}_{\text{mag}}$ is the analogous masked mean-squared error on
the per-frame min--max-normalised magnitude with magnitude weight
$\lambda_{M}=1$. To suppress spurious high-frequency texture in
extraluminal regions, an isotropic total-variation penalty
\cite{rudin1992tv} is applied exclusively to voxel pairs that lie
outside the mask along each Cartesian axis,
\begin{equation}
    \mathcal{L}_{\text{TV}}
    = \sum_{q}\sum_{d\in\{x,y,z\}}
      \frac{\sum_{\mathbf{x}} \lvert
      \hat{q}(\mathbf{x}+\mathbf{e}_d)-\hat{q}(\mathbf{x}) \rvert\,
      \Omega(\mathbf{x},d)}
           {\sum_{\mathbf{x}} \Omega(\mathbf{x},d) + \varepsilon},
\end{equation}
where
$\Omega(\mathbf{x},d) =
 (1-\mathcal{M}(\mathbf{x}))\,(1-\mathcal{M}(\mathbf{x}+\mathbf{e}_d))$
is the joint outside indicator and
$\lambda_{\text{TV}}=10^{-5}$ provides a gentle smoothness prior
without over-regularising the extraluminal signal.

Image-reconstruction quality was assessed by VENC-normalised
mean-squared error (MSE) of the predicted velocity field. The
intraluminal velocity error was computed inside the CoW mask both
as a total
\begin{equation}
    \mathrm{MSE}_{\text{vel}}
    = \frac{\sum_{\mathbf{x}}
            \sum_{q\in\{u,v,w\}}\!\!
            (\hat{q}(\mathbf{x})-q^{*}(\mathbf{x}))^{2}\,
            \mathcal{M}(\mathbf{x})}
           {\sum_{\mathbf{x}}\mathcal{M}(\mathbf{x}) + \varepsilon}
\end{equation}
and per component
\begin{equation}
    \mathrm{MSE}_{q}
    = \frac{\sum_{\mathbf{x}}
            (\hat{q}(\mathbf{x})-q^{*}(\mathbf{x}))^{2}\,
            \mathcal{M}(\mathbf{x})}
           {\sum_{\mathbf{x}}\mathcal{M}(\mathbf{x}) + \varepsilon},
    \qquad q \in \{u,v,w\},
\end{equation}
and an analogous extraluminal quantity --- the background MSE
$\mathrm{MSE}_{\text{bg}}$ --- was obtained by substituting
$\mathcal{M}$ with $1-\mathcal{M}$ in the numerator and denominator.
Both metrics are reported in units of $10^{-3}$ in the
VENC-normalised range.

Vascular-geometry fidelity was evaluated by re-segmenting each
predicted 4D Flow volume with the same ResEnc-M ensemble described
above and comparing the resulting binary CoW mask $P$ to the
subject-specific 7\,T reference mask $R$ in the 3\,T geometry. Volumetric
overlap was quantified by the Dice score,
\begin{equation}
    \mathrm{Dice}(P,R)
    = \frac{2\lvert P \cap R \rvert}{\lvert P \rvert + \lvert R \rvert}.
\end{equation}
Boundary agreement was quantified in physical units from
triangulated surfaces $S_P$ and $S_R$ extracted by marching cubes.
Letting
$d(p,S_R) = \min_{r \in S_R} \lVert p - r \rVert_{2}$ denote the
directed point-to-surface distance, the symmetric mean surface
distance is
\begin{equation}
    \mathrm{sMSD}(P,R)
    = \tfrac{1}{2}\!\left(
        \frac{1}{|S_P|}\sum_{p\in S_P} d(p,S_R)
        + \frac{1}{|S_R|}\sum_{r\in S_R} d(r,S_P)
      \right),
\end{equation}
and the worst-case boundary mismatch is captured by the Hausdorff
distance,
\begin{equation}
    \mathrm{HD}(P,R)
    = \max\!\left(
        \max_{p\in S_P} d(p,S_R),\;
        \max_{r\in S_R} d(r,S_P)
      \right).
\end{equation}
Dice is reported as higher-is-better in $[0,1]$; sMSD and HD are
reported in millimetres and lower-is-better. Together, these metrics
quantify both the volumetric overlap and the boundary fidelity of the
reconstructed vascular geometry that the velocity field implies.

% FIGURE — Loss partition and evaluation regions.
% Conceptual diagram: a representative axial slice with the binary mask
% M partitioning the volume into a fluid (intraluminal) region and an
% extraluminal background; arrows pointing from each region to the
% corresponding loss term (L_vel inside M, L_TV outside M, L_mag inside
% M for the magnitude channel). A second small panel shows the
% prediction-vs-reference surfaces S_P and S_R schematically with arrows
% illustrating sMSD and HD.
\begin{figure}[!htbp]
    \centering
    % Insert figure here
    \caption{Loss partition and evaluation regions. The CoW mask
    $\mathcal{M}$ separates fluid voxels (where
    $\mathcal{L}_{\text{vel}}$ and $\mathcal{L}_{\text{mag}}$ are
    evaluated and intraluminal MSE is reported) from extraluminal voxels
    (where $\mathcal{L}_{\text{TV}}$ and background MSE are computed).
    Geometry metrics are obtained from the predicted and reference
    surfaces $S_P$ and $S_R$.}
    \label{fig:loss}
\end{figure}

\subsection*{Data preparation, training, and inference}

Prior to network input, all velocity components were normalised by the
global VENC,
\begin{equation}
    \tilde{q} = q/\mathrm{VENC}_{\text{global}},
    \qquad
    \mathrm{VENC}_{\text{global}}
    = \max(\mathrm{VENC}_u,\mathrm{VENC}_v,\mathrm{VENC}_w),
\end{equation}
yielding values approximately in $[-1,1]$, and magnitude images were
rescaled to $[0,1]$ by per-frame min--max scaling
\cite{cardoso2022monai}. In both task configurations the reconstruction
target was the registered 7\,T velocity and magnitude volume after
masking by $\mathcal{M}$,
$\mathbf{V}^{*}(\mathbf{x},t)=\mathbf{V}^{\text{reg}}_{7T}(\mathbf{x},t)\,
\mathcal{M}(\mathbf{x})$, confining supervision to the haemodynamically
relevant intraluminal region. CBAM-DNS was trained at native 3\,T
resolution ($r=1$), with $\mathbf{V}^{\text{reg}}_{3T}$ as input and
$\mathbf{V}^{*}$ as target, so that the network simultaneously bridges
the SNR gap between the two field strengths, corrects residual contrast
differences arising from distinct hardware and reconstruction chains,
and recovers finer vascular detail visible only at 7\,T. CBAM-SupRes
was trained with a spatially downsampled input
$\mathbf{V}^{\mathrm{LR}}_{3T}(\mathbf{x},t)
 = \mathcal{Z}_{0.5}[\mathbf{V}^{\text{reg}}_{3T}(\mathbf{x},t)]$,
obtained by isotropic trilinear downsampling, and the same masked
7\,T target, combining $r=2$ upsampling and domain translation in a
single end-to-end pass.

Each task was trained in two complementary input strategies that are
reported side by side in the results. In the \emph{synthetic-noise}
strategy, the network sees the original 3\,T (or downsampled 3\,T)
volume with synthetic non-vascular noise injected into the extraluminal
region during training, and operates without a mask at inference time.
In the \emph{masked-input} strategy, the predicted CoW segmentation is
used to zero the extravascular signal of the 3\,T input prior to the
network for both training and inference; no noise augmentation is then
applied. The synthetic-noise strategy targets mask-free deployment and
emphasises robustness to background patterns, whereas the masked-input
strategy emphasises geometric fidelity by removing the extraluminal
domain shift altogether.

Two augmentation strategies were applied online. Vector-aware geometric
augmentation drew random 90\textdegree{} rotations from the planes
$\{xy,xz,yz\}$ with multiplicities $\{1,2,3\}$ at probability $0.5$,
and the three velocity components were permuted and sign-corrected as
a proper vector field while the mask and magnitude were rotated as
scalar fields. Synthetic non-vascular noise augmentation (used in the
synthetic-noise strategy only) injected per-batch amplitudes drawn
from one of several statistical distributions (zero-mean normal,
Student-$t$, Laplace, skew-normal, or generalised normal) into the
extraluminal region of the input,
\begin{equation}
    \tilde{\mathbf{V}}^{\mathrm{LR}}(\mathbf{x})
    = \mathbf{V}^{\mathrm{LR}}(\mathbf{x})
      + \eta(\mathbf{x})\,(1-\mathcal{M}(\mathbf{x})),
\end{equation}
encouraging mask-focused learning and preventing the network from
overfitting to specific noise patterns in the background.

Training was conducted on three-dimensional patches of size $16^{3}$
voxels in the input space, which corresponded to $32^{3}$ voxels in
the output space for the $\times 2$ configuration
\cite{peiffer2024patch}. For each case, 64 patches were drawn per epoch
with one cardiac frame selected uniformly at random per sample; 16
deterministic centre patches were used for validation, and an optional
minimum-mask-coverage criterion ensured that each sampled patch
contained a sufficient proportion of vascular voxels. Optimisation used
Adam \cite{kingma2015adam} with an initial learning rate of
$2{\times}10^{-4}$ and weight decay $5{\times}10^{-7}$. A
ReduceLROnPlateau scheduler monitored the validation loss with a
patience of ten epochs and halved the learning rate (floor
$10^{-6}$), and an early-stopping criterion with patience twenty
epochs, complemented by an overfitting detector that terminated
training when the training loss kept decreasing while the validation
loss increased for fifteen consecutive epochs, prevented over-fitting.
Models were trained with mini-batches of four samples; the checkpoint
attaining the lowest validation loss was retained for inference.

All results are reported under \emph{leave-one-out cross-validation}
(LOOCV) at the patient level. With $n=7$ subjects, seven independent
folds were trained: in each fold, one subject was held out as the test
case and the remaining six formed the training set; no subject
appeared in both subsets. The metrics reported in the results are
mean $\pm$ standard deviation across the seven held-out subjects.

% FIGURE — Task configurations and LOOCV.
% Two-panel schematic. (a) The two task configurations: x1 maps the
% registered 3T volume to the masked 7T target at the same resolution;
% x2 first downsamples the 3T input by a factor of 0.5 with trilinear
% interpolation and then maps it to the masked 7T target at full
% resolution via the r = 2 network. (b) A 7-row LOOCV chart: in each row
% one subject is highlighted red (test), one yellow (validation), and
% five blue (training).
\begin{figure}[!htbp]
    \centering
    % Insert figure here
    \caption{Task configurations and patient-level leave-one-out
    cross-validation. (a)~$\times 1$ domain translation and $\times 2$
    joint domain translation and super-resolution share the same
    backbone and target. (b)~Seven LOOCV folds; one subject is held out
    as test and the remaining six train the model in every fold.}
    \label{fig:tasks_loocv}
\end{figure}

At inference time, full 4D volumes were reconstructed with an
overlap-aware sliding-window strategy
\cite{cardoso2022monai,ronneberger2015unet}. For each time frame the
volume was partitioned into overlapping $16^{3}$ patches with $25\%$
overlap, and predicted patches were recombined by Gaussian-weighted
blending at overlapping boundaries,
\begin{equation}
    \hat{\mathbf{V}}(\mathbf{x})
    = \frac{\sum_{k} w_{k}(\mathbf{x})\,\hat{\mathbf{V}}_{k}(\mathbf{x})}
           {\sum_{k} w_{k}(\mathbf{x})},
\end{equation}
where $w_k(\mathbf{x})$ is a Gaussian window centred on patch $k$,
eliminating patch-edge discontinuities. The reconstructed volumes were
masked by $\mathcal{M}$ to constrain the output to the CoW region for
the intraluminal metrics, and the predicted velocity components were
finally rescaled to physical units via
$\hat{q}_{\text{phys}} = \hat{q}\cdot\mathrm{VENC}$. For the geometry
evaluation, the reconstructed magnitude (and, when unavailable, the
synthesised angiographic volume of Figure~\ref{fig:segmentation}) was passed
through the ResEnc-M segmentation ensemble to produce the predicted
binary CoW mask $P$ that was subsequently compared against the
subject-specific 7\,T reference $R$ to obtain Dice, sMSD, and HD.

% FIGURE — Inference and qualitative reconstruction.
% Comprehensive figure showing, on a representative held-out subject:
% (a) velocity vector maps (mid-systolic frame) for the 3T input, the 7T
% reference, and the network prediction for the x1 and x2 configurations
% on axial slices through the Circle of Willis; (b) voxel-wise relative
% speed error maps overlaid on the angiographic volume; (c) 3D
% streamline plots of the reconstructed velocity field within the CoW
% mask; (d) the re-segmented predicted CoW surface (S_P) overlaid on the
% reference surface (S_R) for the same subject. Zoomed inset on the MCA
% bifurcation recommended.
\begin{figure}[!htbp]
    \centering
    % Insert figure here
    \caption{Inference and qualitative reconstruction on a
    representative held-out subject. Velocity vector maps, voxel-wise
    error maps, and 3D streamlines for the $\times 1$ and $\times 2$
    configurations, together with the re-segmented predicted CoW
    surface overlaid on the 7\,T reference surface used to compute
    Dice, sMSD, and HD.}
    \label{fig:inference}
\end{figure}

% =====================================================================
%  END OF METHODOLOGY SECTION
% =====================================================================
